\sectionguard
\section*{The Propers: Notable and Quotable}

Four verified afterlives of wording from the scriptural propers, each a use
that moves the phrase into a register its psalmist or sage did not occupy.
Straight exegesis, devotional reuse, and bare quotation are excluded; they
belong in the commentary above. Two of the four rise from the same chant
clause, and that concentration is itself the finding: of everything these
propers contain, it is the plea \latin{exsúrge, Dómine, et iúdica causam
tuam} that history armed --- on both sides of the same war.

\begin{afterlife}{Introit and Gradual, Ps. 73:22 --- the bull that opened Luther's excommunication, its satire, and its bonfire, 1520}
\textbf{Appointed phrase.} \latin{Exsúrge, Dómine, et iúdica causam tuam},
with the Gradual's continuation \latin{memor esto \dots} The chant sings the
liturgical Old-Latin form --- \latin{Dómine} where the Clementine reads
\latin{Deus} --- and it is exactly the liturgical form the bull opens with.

\textbf{Later use and locus.} Three scenes. (1) Leo X's condemnation of
Martin Luther opens: \latin{Exsurge, Domine, et judica causam tuam, memor
esto improperiorum tuorum, eorum, quae ab insipientibus fiunt tota die},
dated 15 June 1520 --- the bull is universally known by the chant's own
incipit, \work{Exsurge Domine}. (2) Ulrich von Hutten republished the bull
with mocking marginal annotations, writing to Luther, ``I send the bull of
[Leo] the Tenth flouted by me in some places.'' (3) Luther to Spalatin,
10 December 1520: ``at nine o'clock, at the eastern gate \dots\ were burned
all the papal books \dots\ and the last bull of Leo X.'' Both letters in
Preserved Smith's translation, \work{Luther's Correspondence}, vol.~1 (1913),
pp.~414--415.

\textbf{The turn.} A congregation's sung plea that God arise and defend his
own cause becomes the juridical opening of an excommunication process --- a
human tribunal speaking the psalm in God's stead --- and its target answers
the appropriation twice over, first with satire in the margins, then with
fire. The verse the Church sings this Sunday as supplication spent the
Reformation as a weapon, and the weapon was contested property.

\textbf{Rights and limit.} The bull's Latin, Hutten's title, and Smith's 1913
English are public domain. The bull quotes the psalm, not the missal chant as
such; what binds it to this Sunday is that chant and bull carry the same
liturgical wording, \latin{Dómine} against the Clementine's \latin{Deus}.
\end{afterlife}

\begin{afterlife}{The same verse --- the standard of the Spanish Inquisition}
\textbf{Appointed phrase.} As above, \latin{Exsúrge, Dómine, et iúdica causam
tuam}.

\textbf{Later use and locus.} The banner of the Spanish Inquisition --- a
cross of rough logs flanked by a sword and an olive branch --- bears the
inscription \textsc{exurge domine et judica causam tuam psalm 73}. Verified
in Bernard Picart's engraving ``Standards and Condemned of the Inquisition,''
in the \work{Histoire générale des cérémonies \dots\ de tous les peuples du
monde}, vol.~2 (Paris: Rollin fils, 1741), held and digitised by the Pitts
Theology Library.

\textbf{The turn.} The plea becomes a processional ensign carried at
autos-da-f\'e, the sword of the sentence on one side of the cross and the
olive of reconciliation on the other --- the psalm's cry for God to judge his
own cause redeployed as an institution's warrant for judging on his behalf.
Set beside the bull above, the same appointed verse armed the tribunal and
named the tribunal's most famous target: prosecution's verse, on both sides
of the divide.

\textbf{Rights and limit.} Engraving and inscription are public domain. The
witness is Picart's 1741 engraving --- a Protestant-published ethnography ---
not a surviving banner; the engraving is nonetheless the standard citable
witness for the standard.
\end{afterlife}

\begin{afterlife}{Communion source-verse, Wis. 16:20 --- Joyce scores a beach scene to the bread of heaven, 1922}
\textbf{Appointed phrase.} \latin{Panem de cælo} with the every-delight
clause; the Communion sings \latin{omne delectaméntum}, ``all that is
delicious.''

\textbf{Later use and locus.} James Joyce, \work{Ulysses} (1922), the
``Nausicaa'' episode: while Gerty MacDowell poses on Sandymount strand and
Bloom watches, the men's retreat at the Star of the Sea church reaches
Benediction --- ``he read out \latin{Panem de coelo praestitisti eis} and Edy
and Cissy were talking about the time all the time'' --- a page after the
choir's \latin{Tantum ergo} has become ``the \latin{Tantumer gosa cramen
tum}'' against the price of Gerty's stockings. Verified in the Project
Gutenberg text of the 1922 edition.

\textbf{The turn.} The versicle drawn from this Sunday's Communion verse ---
bread from heaven, having in it every delight --- is run as liturgical
soundtrack to profane appetite: Joyce lets the manna of every taste and the
window-shopping of every taste interleave sentence by sentence, and the irony
cuts both ways at once.

\textbf{Rights and limit.} \work{Ulysses} (1922) is public domain in the
United States. Joyce quotes the Benediction versicle form
(\latin{præstitísti eis}), not the missal Communion's adapted \latin{dedísti
nobis}; the dependence route is the Benediction rite, which takes the same
Wisdom verse this antiphon takes.
\end{afterlife}

\begin{afterlife}{Offertory, Ps. 30:15--16 --- Browning turns the surrendered times into Victorian optimism, 1864}
\textbf{Appointed phrase.} \latin{In mánibus tuis témpora mea} --- ``my times
are in thy hands.''

\textbf{Later use and locus.} Robert Browning, ``Rabbi Ben Ezra''
(\work{Dramatis Person\ae}, 1864), stanza 1: ``Grow old along with me! / The
best is yet to be, / The last of life, for which the first was made: / Our
times are in His hand / Who saith `A whole I planned, / Youth shows but half;
trust God: see all, nor be afraid!'''

\textbf{The turn.} In the psalm the clause is a hunted man's surrender of his
fate to the God he trusts against enemies and persecutors; in Browning it
becomes the signature of Victorian meliorism about aging --- the surrendered
times now underwrite ``the best is yet to be.'' The verbal point is finer
than it looks: the missal's old psalter sings \latin{témpora mea}, ``times,''
with the Hebrew, where the Clementine and Douay read ``lots'' --- so the
Offertory carries precisely the wording English poetry knows.

\textbf{Rights and limit.} Browning (1864) is public domain. His route is the
English Bible's ``My times are in thy hand,'' not the Latin antiphon; the
agreement of missal and King James against the Clementine is a textual
convergence, disclosed as such, not a chain of dependence.
\end{afterlife}
